\documentclass[a4j,10pt]{jarticle}
\usepackage{amsmath}
\usepackage{amsthm}
\usepackage{amssymb}
\usepackage[dvipdfmx]{graphicx}
\usepackage{url}
\usepackage{bm}

\newtheorem{theorem}{Theorem}[section]
\newtheorem{definition}[theorem]{Definition}
\newtheorem{proposition}[theorem]{Proposition}
\newtheorem{lemma}[theorem]{Lemma}
\newtheorem{corollary}[theorem]{Corollary}

\pagestyle{empty}
\setlength\textheight{255mm}
\setlength\textwidth{170mm}
\setlength\topmargin{-5mm}
\setlength\oddsidemargin{-5mm}
\setlength\headheight{0mm}
\setlength\headsep{0mm}
\setlength\columnsep{8mm}

\newcommand{\tour}{順位戦}

%========================================
\title{
    \LARGE\bf
    素数大富豪\tour \quad 概要
}
\author{はち}
\date{2025-03-13}
%========================================
\begin{document}
    \maketitle
    \thispagestyle{empty}

    \section{開催目的}
    幅広い実力のプレイヤーが参加しやすい大会を開催する．

    \section{会期}
    参加者募集期間：$04$月$01$日から$05$月$31$日まで．

    大会開催期間：$06$月$01$日から翌年$01$月$31$日までの$8$ヶ月．

    \section{参加}
    Google Forms を用いて参加者を募集する．

    普段素数大富豪をするときに使用しているハンドルネームがある場合はそれを用いて参加してください．

    \section{ルール}
    \section{概要}
    \tour 変則リーグ戦．仮順位が近い者同士で対戦し順位を決定する．

    \subsection{仮順位}
    最後に参加した\tour を第$x$期，そのときの順位を$y$位，参加者募集期間開始時のレーティングを$z$とする．
    ただし，
    \begin{itemize}
        \item 過去の\tour に参加したことがない者は，$x=-\infty$, $y=\infty$ とする．
        \item レーティングを持ってない者は，$z=-\infty$ とする．
    \end{itemize}

    これらを用いて仮順位を以下のように決定する．

    \begin{enumerate}
        \item $y-x$ が小さい者を上の順位とする．
        \item 1. が同じ場合，$x$ が大きい者を上の順位とする．
        \item 2. も同じ場合，$z$ が上の者を上の順位とする．
    \end{enumerate}

    \subsection{1試合の流れ}
    
    \subsection{順位決定方法}

\end{document}

%latexmk -latex=platex -pdfdvi rule_01.tex
